% Part of the team paper; included by main.tex (and by nathan_sections.tex for the standalone preview).
\providecommand{\inprogress}[1]{\textcolor{red}{[In progress: #1]}}

\section{Method}
\label{sec:method}

\subsection{Task: Push-T with a hidden goal}
\label{sec:task}
We use \texttt{PushTVarGoal-v1}, a variable-goal version of ManiSkill's \texttt{PushT-v1}~\cite{tao2024maniskill3}
(the Push-T task originates in~\cite{florence2021implicit}), written for this project (Appendix~\ref{sec:ppo}). A Panda arm with a stick
end-effector pushes a T-shaped block on a table to a goal pose. The full state $s_t \in \mathbb{R}^{31}$ contains joint positions and
velocities, the tool-tip pose, the goal position $g$ (3 numbers) and the T's pose. Actions $a_t \in [-1,1]^7$ are joint-position
deltas (\texttt{pd\_joint\_delta\_pos}), and an episode lasts $T = 100$ steps.

\textbf{Goals and quadrants.} Each episode samples a goal position from one of four quadrants $Q_1,\dots,Q_4$ around a fixed
origin: $\pm 10$\,cm in $x$ and $\pm 20$\,cm in $y$, with a 2--3\,cm band around each axis excluded so the quadrants are
separated. The goal orientation is fixed. %CHECK quadrant numbering convention and which axis has the 2 vs 3 cm band
The T starts 8--30\,cm from the goal, and the arm starts from one of 600 precomputed poses.

\textbf{Hiding the goal.} The goal-conditioned teachers (Section~\ref{sec:teachers}) see $s_t$. The transformer only sees
$\phi(s_t) \in \mathbb{R}^{28}$, the state with the goal removed. Within an episode the goal is constant, so the only way the
transformer can learn where to push is from the demonstrations in its prompt.

\textbf{Metrics.} An episode is a success at step $t$ when the T overlaps the goal pose closely enough under ManiSkill's
\texttt{PushT-v1} criterion. %CHECK state the exact threshold
We report \emph{success-once} (success at any step) and \emph{success-at-end} (success at the final step). The gap between them
measures whether a policy that reaches the goal also stops there.

\subsection{Goal-conditioned teachers}
\label{sec:teachers}

\textbf{PPO teacher} (Appendix~\ref{sec:ppo}). A goal-conditioned PPO~\cite{schulman2017ppo} policy trained for 50M
environment steps with ManiSkill's PPO baseline on goals from all four quadrants. We use its deterministic (mean) action.
%CHECK with Jimmy: training details (Appendix sec:ppo)

\textbf{BC teacher, distilled from PPO.} We distill the frozen PPO teacher into an MLP
$\pi_{\mathrm{BC}}: \mathbb{R}^{31} \to [-1,1]^7$ (two hidden layers of 256 units, ReLU, $\tanh$ output). To avoid the compounding
errors of plain behavior cloning~\cite{ross2011dagger}, the training states come from \emph{noisy} teacher rollouts. A Gaussian
perturbation $\epsilon \sim \mathcal{N}(0,\sigma^2 I)$ with $\sigma \in \{0, 0.1, 0.2, 0.3\}$ (cycled across batches) is added to
the executed action, while the label is always the teacher's clean action $\pi_{\mathrm{PPO}}(s_t)$, as in DART~\cite{laskey2017dart}.
We collect 16{,}384 episodes (1.64M transitions) on goals in all four quadrants and minimize the mean squared action error with Adam
(learning rate $3\times10^{-4}$, cosine decay, batch 1{,}024, 200k steps). We do not use DAgger: in an earlier fixed-goal experiment
(Appendix~\ref{app:teacher-history}) it added nothing measurable over noise injection.

\textbf{Teacher quality.} On 100 test episodes per quadrant (seeds fixed per quadrant), the PPO teacher scores 93.8\% / 91.5\% mean
success-once / success-at-end and the BC teacher 91.5\% / 88.8\%. Both reach at least 89\% success-once in every quadrant, including $Q_4$.
A transformer failure on $Q_4$ goals is therefore not explained by weak teachers there.
Both teachers see $Q_4$ goals in training. Only the transformer is restricted to $Q_1$--$Q_3$.

\subsection{Demonstration data}
\label{sec:data}
For each split we draw a fixed set of goal positions from its quadrants, independently and with its own seed
(Table~\ref{tab:splits}). For each goal and teacher we roll out 10 noise-free episodes from different initial states, recording
$(\phi(s_t), a_t)$ for all 100 steps. The goal is recorded only to group episodes by goal and is never given to the model.
Test goals are new positions: none appears in training.

\begin{table}[h]
\centering
\caption{Goal splits for the transformer. Each goal has 10 demonstrations per teacher. Teacher columns give the success-once /
success-at-end of those demonstrations (\%).}
\label{tab:splits}
\begin{tabular}{lccccc}
\toprule
Split & Quadrants & Goals & Demos (both teachers) & PPO demos & BC demos \\
\midrule
train & $Q_1$--$Q_3$ & 2{,}048 & 40{,}960 & 93.8 / 89.3 & 93.1 / 86.8 \\
val & $Q_1$--$Q_3$ & 128 & 2{,}560 & 93.5 / 87.1 & 92.2 / 84.8 \\
test (in distribution) & $Q_1$--$Q_3$ & 128 & 2{,}560 & 94.1 / 91.1 & 93.4 / 87.7 \\
test (held out) & $Q_4$ & 128 & 2{,}560 & 94.5 / 93.4 & 92.7 / 90.5 \\
\bottomrule
\end{tabular}
\end{table}

\subsection{In-context transformer}
\label{sec:icl-model}

\textbf{Prompt.} Following the function-class setup of Garg et al.~\cite{garg2022can}, a prompt contains examples of a hidden
``task'' followed by a query. Here the task is a goal, and the examples are $k$ whole demonstrations of it by one teacher.
Each demonstration $j$ is subsampled every $r = 4$ steps, keeping the last step so the final T pose is included:
$\mathcal{S} = \{T-1, T-1-r, \dots\}$, which gives $P = |\mathcal{S}| = 25$ state--action pairs. With demonstrations
$\tau^1,\dots,\tau^k$ of goal $g$ by teacher $m$, the prompt is
\begin{equation}
  X_k = \big(\phi(s^1_{t}), a^1_{t}\big)_{t \in \mathcal{S}},\ \dots,\ \big(\phi(s^k_{t}), a^k_{t}\big)_{t \in \mathcal{S}},\ \phi(s^q),
  \label{eq:prompt}
\end{equation}
one token per state and per action, so $2kP + 1$ tokens ($401$ at $k = 8$). The query state $s^q$ is a uniformly random step of a
\emph{different} demonstration $\tau^q$ of the same goal and teacher; its label is the teacher's action $a^q$ there.

\textbf{Architecture.} States and actions are embedded by separate linear maps,
$e(s) = W_s\,(\phi(s) - \mu)/\sigma$ and $e(a) = W_a\,a$, where $\mu, \sigma$ are per-dimension training-set statistics. The token
sequence goes through a causal GPT-2 backbone~\cite{radford2019gpt2} with learned absolute position embeddings: 12 layers, width 256,
8 heads and no dropout, 9.75M parameters in total. A linear head reads the last hidden state,
$\hat a = W_o\, h_{2kP+1}$. Unlike Garg et al., who predict at every query position, we predict only at the final token.

\textbf{Training objective.} Each step samples one context size $k$ uniformly from $\{0,1,2,4,8\}$ and then, for each of 64
prompts, a teacher $m$ and a training goal $g$ uniformly. Within a goal, demonstrations are ordered successful-first (random
otherwise), the first $k$ form the support and the next one supplies the query. We minimize
\begin{equation}
  \mathcal{L}(\theta) = \mathbb{E}_{k,\,m,\,g,\,\tau^{1:k},\,\tau^q,\,t}\,
  \big\| f_\theta(X_k) - a^q_t \big\|_2^2 .
  \label{eq:loss}
\end{equation}
At $k = 0$ the prompt is only $\phi(s^q)$, so the model must output a goal-independent action. That gives the no-context baseline.
Table~\ref{tab:hparams} lists the hyperparameters. The architecture, learning rate and batch size follow Garg et al.~\cite{garg2022can}
and were not tuned. %CHECK: confirm Garg et al. use 12 layers / 256 / 8 heads, lr 1e-4, batch 64

\begin{table}[h]
\centering
\caption{Transformer hyperparameters. Nothing was tuned on test data. The only earlier run is a 3k-step pilot with the same settings.}
\label{tab:hparams}
\small
\begin{tabular}{p{0.19\linewidth}p{0.37\linewidth}p{0.33\linewidth}}
\toprule
Setting & Value & Reason \\
\midrule
Layers / width / heads & 12 / 256 / 8 & GPT-2 setup of~\cite{garg2022can} \\
Optimizer & AdamW~\cite{loshchilov2019adamw}, lr $10^{-4}$, no weight decay & lr as in~\cite{garg2022can} \\
Batch size & 64 prompts & as in~\cite{garg2022can} \\
Schedule & 1k-step linear warmup, cosine~\cite{loshchilov2017sgdr} to 10\% & standard; not tuned \\
Training steps & 100k & compute budget (3.6\,h on one L4); val.\ loss still decreasing \\
Context sizes $k$ & $\{0,1,2,4,8\}$, uniform & no context up to 8 demos \\
Subsampling stride $r$ & 4 (25 pairs per demo) & $k = 8$ fits in 1{,}024 positions \\
Precision & bf16 autocast, gradient clip 1.0 & speed \\
Checkpoint selection & lowest val.\ MSE averaged over $k$ (correct goal) & selected the final step \\
\bottomrule
\end{tabular}
\end{table}

\subsection{Evaluation}
\label{sec:eval}
All evaluation uses the frozen checkpoint and goals the transformer never trained on.

\textbf{Controls.} The design has to show that success comes from reading the context. We therefore evaluate every setting under
three conditions:
\begin{itemize}
  \item \emph{correct goal}: demonstrations of the evaluation goal;
  \item \emph{wrong goal}: demonstrations of a different goal from the same split, with everything else identical;
  \item \emph{no demos}: $k = 0$.
\end{itemize}
A model that ignores its context scores the same under all three.

\textbf{Teacher-forced action error.} For each split and $k$, the mean squared error of $\hat a$ on 1{,}024 sampled prompts.
The correct- and wrong-goal conditions use the same queries.

\textbf{Closed-loop success.} For each split we take 128 goals and 4 new initial states per goal, 512 episodes per condition. The
support for goal $g$ is fixed: the first $k$ demonstrations of $g$ (successful first) collected by the context teacher. Every
condition, including the two goal-seeing teachers, runs from the \emph{same} 512 initial states, so differences between conditions
are paired. Wrong-goal episodes use the support of goal $(g + 64) \bmod 128$.

\begin{algorithm}[h]
\caption{Closed-loop evaluation of one condition (vectorized over the 512 episodes in practice)}
\label{alg:closed-loop}
\begin{algorithmic}[1]
\Require frozen $f_\theta$; support $D$ of $k$ subsampled demonstrations (Eq.~\ref{eq:prompt}); goal $g$; initial-state seed
\State $s_0 \gets \textsc{Reset}(g, \text{seed})$ \Comment{same seed for every condition}
\For{$t = 0, \dots, T-1$}
  \State $a_t \gets \mathrm{clip}\big(f_\theta(D, \phi(s_t)),\, -1,\, 1\big)$ \Comment{the goal is never an input}
  \State $s_{t+1} \gets \textsc{Step}(s_t, a_t)$; record whether the episode is a success at $t+1$
\EndFor
\State \Return success-once, success-at-end
\end{algorithmic}
\end{algorithm}

\textbf{Uncertainty.} Success rates come with 95\% Wilson score intervals~\cite{wilson1927probable} for a binomial proportion over
$n = 512$ episodes: about $\pm 3$ points near 90\% and $\pm 4$ near 65\%. These intervals treat episodes as independent. They ignore
that the 4 episodes of a goal share it, and they do not capture variation across training seeds.
\inprogress{only one training seed so far; seeds 1--2 will give between-seed error bars.}
\inprogress{a goal-estimate baseline: estimate $g$ from the final T pose in the demonstrations and give it to the BC teacher.}

